\documentclass[11pt]{amsart}

\usepackage[T1]{fontenc}
\usepackage{lmodern}
\usepackage{microtype}
\usepackage{mathtools,amssymb,amsthm}
\usepackage[hidelinks]{hyperref}

\hypersetup{
  pdftitle={An explicit OF(6,4) and an explicit OF(8,4)}
}

\newtheorem{theorem}{Theorem}
\newtheorem{lemma}[theorem]{Lemma}
\newtheorem{proposition}[theorem]{Proposition}
\newtheorem{corollary}[theorem]{Corollary}
\theoremstyle{remark}
\newtheorem*{remark}{Remark}

\newcommand{\OF}{\mathrm{OF}}
\newcommand{\Km}[2]{K_{#1\times #2}}
\newcommand{\Zp}{\mathbb Z}

\title[An explicit $\OF(6,4)$ and $\OF(8,4)$]
{An explicit $\OF(6,4)$ and an explicit $\OF(8,4)$}

\date{}

\subjclass[2020]{05C70, 94B25, 05C25}
\keywords{one-factorization, complete multipartite graph, distance constraint,
constant-weight code, prescribed automorphism group, starter}

\begin{document}

\begin{abstract}
An $\OF(n,g)$ is a one-factorization of the complete multipartite graph
$\Km ng$ in which no one-factor contains two edges joining the same pair of
parts; equivalently, an optimal decomposition of $\Km ng$ with respect to
distance $3$. Tan, Zhou and Etzion state that for even $g\ge4$ an $\OF(n,g)$
exists if and only if $n>g$, \emph{possibly except} odd $n$ with
$g+3\le n\le 2g-3$ and even $n$ with $g+2\le n\le 4g-4$; at $g=4$ that
exception list is $\{6,8,10,12\}$, and they settle $n=10$. We exhibit a
$12$-edge base factor whose $20$ translates under $\Zp_5\times\Zp_4$ form an
$\OF(6,4)$, and a $16$-edge base factor whose $28$ translates under
$\Zp_7\times\Zp_4$ form an $\OF(8,4)$. Both base factors are printed in full,
and verifying each is a finite inspection of its twelve resp.\ sixteen edges
together with three short lemmas. Granting the statements we quote from Tan,
Zhou and Etzion, their doubling lemma turns the $\OF(6,4)$ into an $\OF(12,4)$
--- which we neither reprove nor write out --- and an $\OF(n,4)$ then exists for
every $n>4$. We do not determine whether an $\OF(6,4)$ or an $\OF(8,4)$ already
follows from known amalgamation--detachment results.
\end{abstract}

\maketitle

\section{Definitions and the quoted exception list}

For integers $n\ge2$ and $g\ge1$ let $\Km ng$ be the complete multipartite
graph with $n$ parts of size $g$, so
$|V(\Km ng)|=ng$ and $|E(\Km ng)|=g^2\binom n2$. Two disjoint edges of
$\Km ng$ are \emph{parallel} if their ends lie in the same pair of parts.
Following Tan, Zhou and Etzion \cite{TZE}, an $\OF(n,g)$ is a
one-factorization of $\Km ng$ in which no one-factor contains a pair of
parallel edges; equivalently, an optimal decomposition of $\Km ng$ with respect
to distance $3$, since a one-factor read as a set of words of length $n$ and
weight $2$ over $\{0,1,\dots,g\}$ has minimum Hamming distance $\ge3$ exactly
when it has no two edges on one pair of parts. The notation carries $g>1$ and
$n>g$. The necessary conditions recorded in \cite{TZE} are that $ng$ be
even and $n>g$; at $g=4$ the parity clause is vacuous, so $n>4$ is the whole of
it.

Part (2) of the main existence theorem of \cite{TZE} reads: for even $g\ge4$
there exists an $\OF(n,g)$ if and only if $n>g$, possibly except odd $n$ with
$g+3\le n\le2g-3$ and even $n$ with $g+2\le n\le4g-4$. This clause and the other
clauses of \cite{TZE} used below are quoted from the e-print cited in the
bibliography, not from a published version. At $g=4$ the
odd range $[g+3,2g-3]=[7,5]$ is empty and the even range $[g+2,4g-4]=[6,12]$
contributes $\{6,8,10,12\}$, the four values that quoted statement leaves open
at $g=4$. Of these, $n=10$ is settled inside \cite{TZE} itself, by their
$\OF(g+1,g)$ for even $g$ (which gives an $\OF(5,4)$) followed by their
doubling lemma. We exhibit objects for $n=6$ and $n=8$.

\begin{theorem}\label{thm:main}
An $\OF(6,4)$ and an $\OF(8,4)$ exist. Both are exhibited in
Section~\textup{\ref{sec:objects}} as the orbit of a single one-factor under a
group of order $20$ resp.\ $28$ acting freely on the edge set.
\end{theorem}

\begin{corollary}\label{cor:spectrum}
Grant the clauses quoted from \cite{TZE} above and the doubling lemma of
\cite{TZE} recorded as Lemma~\textup{\ref{lem:double}} below. Then an $\OF(n,4)$
exists if and only if $n>4$.
\end{corollary}

Only Theorem~\ref{thm:main} is established here. The rest of
Corollary~\ref{cor:spectrum} is quoted: $n=12$ is Theorem~\ref{thm:main} put
through the doubling lemma of \cite{TZE}, no $\OF(12,4)$ being written out
below, and every other $n>4$ --- including $n=10$ --- is already in \cite{TZE}.

We claim no priority for the two objects. We have not determined whether an
amalgamation--detachment theorem controlling per-colour degrees and total
multiplicities at once, such as that of Bahmanian and Rodger
\cite{BahmanianRodger}, yields these cells directly, and we have not established
that no $\OF(6,4)$ or $\OF(8,4)$ appears in the literature. What is offered is
two explicit objects with a verification a reader can carry out by hand.

\section{The two objects}\label{sec:objects}

Fix an odd integer $p\ge3$ and an integer $g\ge1$, and put $n=p+1$. Take the
parts of $\Km ng$ to be indexed by $\Zp_p\cup\{\infty\}$ and the $g$ vertices
of a part by $\Zp_g$, so that
\[
 V=\bigl(\Zp_p\cup\{\infty\}\bigr)\times\Zp_g,
 \qquad
 P_u=\{u\}\times\Zp_g \quad(u\in\Zp_p\cup\{\infty\})
\]
and $E=E(\Km ng)$ is the set of pairs of vertices lying in different parts. Let
\[
 G=\Zp_p\times\Zp_g,
 \qquad
 (t,s)\colon (x,a)\longmapsto (x+t,\;a+s),
 \qquad \infty+t=\infty ,
\]
a group of order $pg=(n-1)g$ acting on $V$, permuting the parts, and therefore
acting on $E$.

\subsection*{An $\OF(6,4)$}
Take $p=5$, $g=4$: six parts, $24$ vertices, $|E|=\binom{24}2-6\binom42=240$,
and $|G|=20$. Let $M$ be the following one-factor of $\Km64$. Each line is one
edge, written as \texttt{part value part value}, with \texttt{inf} for
$\infty$ and values in $\Zp_4$:
\begin{center}
\begin{minipage}{0.42\textwidth}
\begin{verbatim}
0 0  1 1
1 2  2 2
3 3  4 1
4 3  0 2
0 1  2 1
1 3  3 1
2 3  4 2
3 2  0 3
inf 0  1 0
inf 1  2 0
inf 2  3 0
inf 3  4 0
\end{verbatim}
\end{minipage}
\end{center}
The $12$ edges cover the $24$ vertices once each, so $M$ is a perfect matching,
and each part contributes the four values $0,1,2,3$.

\subsection*{An $\OF(8,4)$}
Take $p=7$, $g=4$: eight parts, $32$ vertices,
$|E|=\binom{32}2-8\binom42=496-48=448$, and $|G|=28$. Let $M_8$ be
\begin{center}
\begin{minipage}{0.42\textwidth}
\begin{verbatim}
0 0  1 0
0 1  2 1
0 2  4 1
0 3  5 0
1 1  2 2
1 2  inf 2
1 3  5 1
2 0  3 2
2 3  6 0
3 0  4 3
3 1  5 2
3 3  6 3
4 0  6 2
4 2  inf 1
5 3  inf 0
6 1  inf 3
\end{verbatim}
\end{minipage}
\end{center}
Again the $16$ edges cover the $32$ vertices once each and each part
contributes $0,1,2,3$.

\begin{theorem}\label{thm:objects}
The $20$ translates $M+\gamma$, $\gamma\in\Zp_5\times\Zp_4$, form an
$\OF(6,4)$, and the $28$ translates $M_8+\gamma$,
$\gamma\in\Zp_7\times\Zp_4$, form an $\OF(8,4)$.
\end{theorem}

\section{The argument}

Theorem~\ref{thm:objects} follows from two lemmas that hold for every odd $p$
and every $g$, plus a finite check on the printed edges.

\begin{lemma}\label{lem:free}
$G$ acts freely on $E$, and $E$ splits into exactly $gn/2$ orbits, each of
size $|G|=(n-1)g$.
\end{lemma}

\begin{proof}
Let $\gamma=(t,s)\in G$ fix an edge $e\in E$.

If $e=\{(x,a),(y,b)\}$ with $x,y\in\Zp_p$, $x\ne y$, then either $\gamma$
fixes each end or it interchanges them. In the second case $x+t=y$ and
$y+t=x$, so $2t\equiv0\pmod p$; as $p$ is odd, $2$ is invertible modulo $p$,
whence $t=0$ and then $x=y$, a contradiction. In the first case $x+t=x$ gives
$t=0$ and $a+s=a$ gives $s=0$, i.e.\ $\gamma$ is the identity.

If $e=\{(\infty,a),(y,b)\}$ with $y\in\Zp_p$, the two ends lie in different
parts and $\infty$ is fixed by every element of $G$, so $\gamma$ cannot
interchange them; fixing them individually gives $t=0$ and $s=0$ as before.

So the stabiliser of every edge is trivial and every orbit has exactly $|G|$
elements. Consequently the number of orbits is
\[
 \frac{|E|}{|G|}
 =\frac{g^2\binom n2}{(n-1)g}
 =\frac{gn}2 . \qedhere
\]
\end{proof}

Since a one-factor of $\Km ng$ has $ng/2$ edges, Lemma~\ref{lem:free} says that
the number of $G$-orbits on $E$ equals the size of a one-factor; so a
one-factor meeting every orbit at most once meets every orbit exactly once.

\begin{lemma}\label{lem:starter}
Let $M\subseteq E$ be a perfect matching of $\Km ng$ such that
\begin{enumerate}
\item[\textup{(T)}] $M$ contains exactly one edge from each $G$-orbit on $E$,
and
\item[\textup{(P)}] the pairs of parts spanned by the edges of $M$ are pairwise
distinct.
\end{enumerate}
Then $\{M+\gamma:\gamma\in G\}$ is an $\OF(n,g)$.
\end{lemma}

\begin{proof}
Each element of $G$ permutes $V$ and maps parts onto parts, so every
$M+\gamma$ is again a perfect matching of $\Km ng$ and, by (P), again has
pairwise distinct part-pairs: no $M+\gamma$ contains parallel edges.

The translates are pairwise edge-disjoint. Suppose
$e_1+\gamma_1=e_2+\gamma_2$ with $e_1,e_2\in M$ and
$\gamma_1,\gamma_2\in G$. Then $e_1$ and $e_2$ lie in one orbit, so $e_1=e_2$
by (T), and then $e_1+\gamma_1\gamma_2^{-1}=e_1$ forces
$\gamma_1=\gamma_2$ by Lemma~\ref{lem:free}. In particular the $|G|$
translates are pairwise distinct, and they carry
\[
 |G|\cdot\frac{ng}2
 =(n-1)g\cdot\frac{ng}2
 =g^2\binom n2
 =|E|
\]
distinct edges, i.e.\ they partition $E$. A partition of $E(\Km ng)$ into
perfect matchings is a one-factorization, and none of its factors has parallel
edges.
\end{proof}

It remains to verify (T) and (P) for the two printed matchings, and for (T) we
need to name the orbits. For $e\in E$ set
\begin{equation}\label{eq:inv}
 \iota(e)=
 \begin{cases}
 (d,\,b-a) & e=\{(x,a),(x+d,b)\},\ x\in\Zp_p,\ 1\le d\le\tfrac{p-1}2,\\[2pt]
 (\infty,\,b-a) & e=\{(\infty,a),(y,b)\},\ y\in\Zp_p ,
 \end{cases}
\end{equation}
the first line being well defined because exactly one of the two orientations
of a finite--finite edge has its part-difference in $\{1,\dots,(p-1)/2\}$.

\begin{lemma}\label{lem:inv}
$\iota$ is a complete invariant of the $G$-orbits on $E$: two edges lie in one
orbit if and only if they have the same value of $\iota$.
\end{lemma}

\begin{proof}
$\iota$ is $G$-invariant, since $(t,s)$ shifts the two part-indices of an edge
by the same $t$ and its two values by the same $s$. The fibre of $\iota$ over
$(d,j)$ with $1\le d\le(p-1)/2$ is
$\{\{(x,a),(x+d,a+j)\}: x\in\Zp_p,\ a\in\Zp_g\}$ and the fibre over
$(\infty,j)$ is $\{\{(\infty,a),(y,a+j)\}: y\in\Zp_p,\ a\in\Zp_g\}$; each has
exactly $pg=|G|$ elements. So $\iota$ has
$g\bigl(\frac{p-1}2+1\bigr)=\frac{g(p+1)}2=\frac{gn}2$ fibres, each of size
$|G|$, and each fibre is a union of orbits. By Lemma~\ref{lem:free} there are
$gn/2$ orbits and each has size $|G|$, so each fibre is a single orbit.
\end{proof}

\begin{proof}[Proof of Theorem~\ref{thm:objects}]
Both matchings were seen to be perfect matchings in Section~\ref{sec:objects},
so by Lemmas~\ref{lem:starter} and~\ref{lem:inv} it suffices to read off
$\iota$ and the part-pairs.

For $M$ at $p=5$ there are $12$ orbits, indexed by
$d\in\{1,2\}$ and $\infty$, each with four values of $b-a$ in $\Zp_4$. In the
order the edges are printed:
\[
\begin{array}{lcccc}
 d=1: & (0,0)(1,1) & (1,2)(2,2) & (3,3)(4,1) & (4,3)(0,2)\\
      & b-a=1 & 0 & 2 & 3\\[4pt]
 d=2: & (0,1)(2,1) & (1,3)(3,1) & (2,3)(4,2) & (3,2)(0,3)\\
      & b-a=0 & 2 & 3 & 1\\[4pt]
 \infty: & (\infty,0)(1,0) & (\infty,1)(2,0) & (\infty,2)(3,0) &
 (\infty,3)(4,0)\\
      & b-a=0 & 3 & 2 & 1
\end{array}
\]
(the two entries in the fourth column of the first two rows use the
orientations $4\mapsto0$ and $3\mapsto0$, of difference $1$ and $2$). Each of
the three rows shows each element of $\Zp_4$ exactly once, so all $12$ values
of $\iota$ are distinct and (T) holds. The part-pairs of $M$ are
\[
\begin{gathered}
 \{0,1\},\ \{1,2\},\ \{3,4\},\ \{4,0\},\ \{0,2\},\ \{1,3\},\\
 \{2,4\},\ \{3,0\},\ \{\infty,1\},\ \{\infty,2\},\ \{\infty,3\},\
 \{\infty,4\},
\end{gathered}
\]
twelve distinct pairs out of $\binom62=15$, so (P) holds.

For $M_8$ at $p=7$ there are $16$ orbits, indexed by $d\in\{1,2,3\}$ and
$\infty$. Reading \eqref{eq:inv} off the printed edges gives, grouped by
index,
\[
\begin{array}{@{}l@{\,}c@{\,}c@{\,}c@{\,}c@{}}
 d=1: & (0,0)(1,0)\mapsto0, & (1,1)(2,2)\mapsto1,
      & (2,0)(3,2)\mapsto2, & (3,0)(4,3)\mapsto3;\\[2pt]
 d=2: & (0,1)(2,1)\mapsto0, & (3,1)(5,2)\mapsto1,
      & (4,0)(6,2)\mapsto2, & (0,3)(5,0)\mapsto3;\\[2pt]
 d=3: & (3,3)(6,3)\mapsto0, & (0,2)(4,1)\mapsto1,
      & (1,3)(5,1)\mapsto2, & (2,3)(6,0)\mapsto3;\\[2pt]
 \infty: & (1,2)(\infty,2)\mapsto0, & (4,2)(\infty,1)\mapsto1,
      & (6,1)(\infty,3)\mapsto2, & (5,3)(\infty,0)\mapsto3 .
\end{array}
\]
Four of the sixteen edges are read in the orientation that starts at the larger
part label, that being the one whose difference lies in $\{1,2,3\}$:
$(0,2)(4,1)$ from $4$ to $0$, since $4+3=0$ in $\Zp_7$, giving $b-a=2-1=1$;
$(0,3)(5,0)$ from $5$ to $0$, giving $3-0=3$; $(1,3)(5,1)$ from $5$ to $1$,
giving $3-1=2$; and $(2,3)(6,0)$ from $6$ to $2$, giving $3-0=3$. Each of the
four rows shows each element of
$\Zp_4$ exactly once, so (T) holds. The part-pairs of $M_8$ are the sixteen
\[
\begin{gathered}
 \{0,1\},\ \{0,2\},\ \{0,4\},\ \{0,5\},\ \{1,2\},\ \{1,5\},\
 \{1,\infty\},\ \{2,3\},\\
 \{2,6\},\ \{3,4\},\ \{3,5\},\ \{3,6\},\ \{4,6\},\ \{4,\infty\},\
 \{5,\infty\},\ \{6,\infty\},
\end{gathered}
\]
distinct, out of $\binom82=28$, so (P) holds.
\end{proof}

\section{What follows from the quoted clauses}\label{sec:spectrum}

\begin{lemma}[\cite{TZE}, doubling lemma; quoted, not proved here]\label{lem:double}
Let $n,g$ be positive integers with $ng$ even and $n>g$. If there exists an
$\OF(n,g)$, then there exists an $\OF(2n,g)$.
\end{lemma}

At $(n,g)=(6,4)$ we have $ng=24$ even and $n=6>4=g$, which is all
Lemma~\ref{lem:double} as quoted asks, so Theorem~\ref{thm:main} and that lemma
give an $\OF(12,4)$; we do not write the object out. (The proof of
Lemma~\ref{lem:double} in \cite{TZE} uses a pair of orthogonal Latin squares of
order $g$ and excludes $g=2,6$; both conditions are satisfied at $g=4$. We have
not reproved the lemma.)

\begin{proof}[Proof of Corollary~\ref{cor:spectrum}]
Necessity is the necessary condition of \cite{TZE} quoted in Section~1:
$n>g=4$. For sufficiency, let $n>4$. By part (2) of the existence theorem of
\cite{TZE}, an $\OF(n,4)$ exists for every $n>4$ outside the exception list
$\{6,8,10,12\}$. Within it, $n=6$ and $n=8$ are Theorem~\ref{thm:main};
$n=12$ is Theorem~\ref{thm:main} with Lemma~\ref{lem:double}; and $n=10$ is
the $\OF(5,4)$ of \cite{TZE} with Lemma~\ref{lem:double}, whose hypotheses hold
at $(n,g)=(5,4)$ since $ng=20$ is even and $5>4$. Every step but the two cases
of Theorem~\ref{thm:main} is a clause of \cite{TZE} taken on trust.
\end{proof}

\section{Verification}

Theorem~\ref{thm:objects} can be checked by hand from the two printed matchings:
Lemmas~\ref{lem:free}--\ref{lem:inv} are three short arguments, and the two
finite checks are the $\iota$-tables and part-pair lists displayed in its proof.

The accompanying program \texttt{verify.py} also expands the factorizations. It
parses the two blocks of Section~\ref{sec:objects} as printed, rebuilds $\Km64$,
$\Km84$ and the two groups, verifies freeness, computes the orbits explicitly
and confirms that $\iota$ is a complete invariant, and expands the $20$ and $28$
translates, checking that each is a perfect matching with pairwise distinct
part-pairs and that the translates partition the $240$ resp.\ $448$ edges. Its
recorded output (\texttt{verify.output.txt}) reports

\begin{center}
\texttt{VERDICT: ALL 56 CHECKS PASS}
\end{center}

\noindent
and states what it does not cover: it does not reprove Lemma~\ref{lem:double}
and constructs no $\OF(12,4)$; it takes the clauses of \cite{TZE} used in
Corollary~\ref{cor:spectrum} on trust, and checks their joint coverage only over
the finite range $5\le n\le4000$; it does not check the wording quoted from
\cite{TZE} against any other version of that source; and it says nothing about
minimality, uniqueness, priority or $g\ge5$.

\begin{thebibliography}{9}

\bibitem{BahmanianRodger}
M.~A. Bahmanian and C.~A. Rodger,
\emph{Multiply balanced edge colorings of multigraphs},
J. Graph Theory \textbf{70} (2012), no.~3, 297--317.
\href{https://doi.org/10.1002/jgt.20617}{doi:10.1002/jgt.20617}.

\bibitem{TZE}
Y.~Tan, J.~Zhou, and T.~Etzion,
\emph{One-factorizations of complete multipartite graphs with distance
constraints},
\href{https://arxiv.org/abs/2602.16319v1}{arXiv:2602.16319v1}, 18 February
2026. Statements are cited descriptively from the \LaTeX{} source of that
e-print; they have not been checked against a published version.

\end{thebibliography}

\end{document}
